\chapter*{Resumen}
\addcontentsline{toc}{chapter}{Resumen}

Esta tesis doctoral contrasta empíricamente las predicciones de un modelo de agencia multitarea sobre los efectos asimétricos de la descentralización hospitalaria en la eficiencia técnica. El modelo teórico, de tres agentes (pagador público, hospital y médico), predice que la descentralización contractual mejora la eficiencia en cantidad de prestaciones pero tiene efectos ambiguos sobre la intensidad de tratamiento, dependiendo del grado de sustituibilidad entre dimensiones del esfuerzo médico (parámetro $\psi_{12}$).

El análisis empírico utiliza métodos de frontera estocástica (SFA) con cuatro diseños complementarios y datos de panel del sistema hospitalario español (SIAE y CNH) para el período 2010--2023. Las predicciones contrastadas son: P1 (la descentralización mejora la eficiencia en cantidad), P2 (la descentralización aumenta la ineficiencia en intensidad), P3 (los hospitales privados de mercado son más eficientes en cantidad que los concertados), P4 (el efecto diferencial de la descentralización es mayor en servicios quirúrgicos que en médicos), y P* (las ineficiencias en ambas dimensiones no son independientes).

Los resultados principales son los siguientes. P1 se confirma robustamente en las siete especificaciones estimadas ($\hat{\alpha}_1 < 0$, significativo al 1\%). P2 no se confirma en su formulación estricta de signo ($\hat{\delta}_1 = -1.477$, signo opuesto al predicho), pero la asimetría entre dimensiones ---la descentralización mejora más la eficiencia en cantidad que en intensidad--- es estadísticamente significativa ($z = -2.09$, $p = 0.037$) y consistente con el modelo bajo sustitución moderada. P3 y P4 se confirman en todos los diseños. P* se confirma con un estadístico de razón de verosimilitudes superior a 10.000.

La principal contribución es la primera contrastación empírica directa de un modelo de agencia multitarea en el sector hospitalario español, con estimación simultánea de fronteras para cantidad e intensidad y un marco teórico que genera predicciones diferenciadas según la estructura organizativa del hospital.

\medskip
\noindent\textbf{Palabras clave:} eficiencia hospitalaria, frontera estocástica, teoría de agencia, multitarea, estructura organizativa, sistema hospitalario español

\medskip
\noindent Clasificación JEL: I11, I18, D82, D86, C23

\chapter*{Abstract}
\addcontentsline{toc}{chapter}{Abstract}

This doctoral thesis empirically tests the predictions of a multitask agency model on the asymmetric effects of hospital decentralization on technical efficiency. The theoretical model, with three agents (public payer, hospital, and physician), predicts that contractual decentralization improves efficiency in the quantity of services but has ambiguous effects on treatment intensity, depending on the degree of substitutability between dimensions of physician effort (parameter $\psi_{12}$).

The empirical analysis uses stochastic frontier analysis (SFA) with four complementary designs and panel data from the Spanish hospital system (SIAE and CNH) for the period 2010--2023. The hypotheses tested are: P1 (decentralization improves quantity efficiency), P2 (decentralization increases inefficiency in intensity), P3 (private market hospitals are more quantity-efficient than contracted hospitals), P4 (the differential effect of decentralization is greater in surgical than in medical services), and P* (inefficiencies in both dimensions are not independent).

The main results are as follows. P1 is robustly confirmed in all seven estimated specifications ($\hat{\alpha}_1 < 0$, significant at 1\%). P2 is not confirmed in its strict sign formulation ($\hat{\delta}_1 = -1.477$, sign opposite to predicted), but the predicted asymmetry between dimensions ---decentralization improves quantity efficiency more than intensity efficiency--- is statistically significant ($z = -2.09$, $p = 0.037$) and consistent with the model under moderate substitution. P3 and P4 are confirmed across all designs. P* is confirmed with a likelihood ratio statistic exceeding 10,000.

The main contribution is the first direct empirical test of a multitask agency model in the Spanish hospital sector, with simultaneous estimation of frontiers for quantity and intensity and a theoretical framework that generates differentiated predictions according to the organizational structure of the hospital.

\medskip
\noindent\textbf{Keywords:} hospital efficiency, stochastic frontier analysis, agency theory, multitasking, organizational structure, Spanish hospital system

\medskip
\noindent JEL Classification: I11, I18, D82, D86, C23

\chapter*{Agradecimientos}
\addcontentsline{toc}{chapter}{Agradecimientos}
% [Pendiente de redacción]
\vspace{2cm}
\begin{center}
\textit{[Pendiente de redacción]}
\end{center}

\tableofcontents
\listoftables
\listoffigures
